% =========================================================================
% Standalone LaTeX Article: Consolidated Priority Matrices v2
% =========================================================================

\documentclass[10pt]{article}

% ---- Essential Packages ----
\usepackage{cite}
\usepackage{amsmath,amssymb,amsfonts}
\usepackage{graphicx}
\usepackage{tikz}
\usetikzlibrary{positioning,shapes.geometric}
\usepackage{hyperref}
\usepackage{enumitem}
\usepackage{booktabs}
\usepackage{longtable}
\usepackage{pdflscape}
\usepackage{geometry}
\usepackage[table]{xcolor}
\usepackage{array}
\usepackage{multirow}

% ---- Page Geometry ----
\geometry{margin=1in}

% ---- Title and Author ----
\title{Consolidated Priority Matrices v2:\\Cybersecurity \& Cyber-Resilience Strategy for the Sovereign Cloud--Edge--IoT Continuum}
\author{Iraklis Symeonidis}
\date{\today}

% ---- Helper macro for gate labels ----
\newcommand{\gatelbl}[1]{#1}

% ---- Helper macro for SSL badges ----
\definecolor{sslredbg}{RGB}{220,20,60}
\definecolor{sslredtx}{RGB}{255,255,255}
\definecolor{sslamberbg}{RGB}{255,165,0}
\definecolor{sslambtx}{RGB}{255,255,255}
\definecolor{ssltealbg}{RGB}{0,128,128}
\definecolor{ssltealtx}{RGB}{255,255,255}

\newcommand{\sslbadge}[3]{%
  \parbox{\linewidth}{\raggedright
    \colorbox{#1}{\scriptsize\parbox{0.92\linewidth}{\raggedright
      \textbf{\textcolor{#2}{#3}}}}}}

\newcommand{\kpistrip}[1]{%
  \par\vspace{2pt}\noindent\rule{\linewidth}{0.2pt}%
  \par\noindent\textit{\scriptsize #1}}

\newcommand{\cellhead}[1]{\textbf{\small #1}\par\smallskip}

% ---- Document Begin ----
\begin{document}

\maketitle

\begin{abstract}
This document presents the consolidated priority matrices for the European Union's cybersecurity and cyber-resilience strategy across the Cloud--Edge--IoT continuum. The matrices organize ten strategic priorities across three clusters (Sovereign Trusted Infrastructure, Cognitive Security and Autonomous Cyber Defence, and Federated Operational Sovereignty) and track their progression through three phases (2025--2030, 2030--2035, 2035--2040). Each priority is assessed across four analytical dimensions: (1) R\&D Maturity, (2) Policy/Regulatory/Standards, (3) Operational Deployment \& Adoption, and (4) Market Uptake \& Strategic Sovereignty.
\end{abstract}

\tableofcontents

\newpage

\section{Introduction}

This document contains three priority progression matrices structured to align with the generalized methodology framework. Each matrix analyzes strategic priorities across three phases and four analytical dimensions, with critical enabler (gate condition) rows indicating the necessary conditions for phase transitions.

~\\

% =========================================================================
% CLUSTER 1 MATRIX
% =========================================================================

\section{Priority Cluster 1: Sovereign Trusted Infrastructure}

% ---- Shared colour definitions --------
\definecolor{dimRI}{RGB}{237,245,253}       % R&I rows    — blue tint
\definecolor{dimPol}{RGB}{237,249,244}      % Policy rows — green tint
\definecolor{dimOps}{RGB}{253,243,237}      % Ops rows    — orange tint
\definecolor{dimMkt}{RGB}{245,244,254}      % Market rows — purple tint
\definecolor{dimGate}{RGB}{240,240,238}     % Gate rows   — neutral grey
\definecolor{dimSub}{RGB}{248,248,246}      % Sub-priority rows
\definecolor{dimHdr}{RGB}{220,220,218}      % Section separator rows
\definecolor{clr1hdr}{RGB}{46,117,182}      % Cluster 1 header bar
\definecolor{clr2hdr}{RGB}{84,130,53}       % Cluster 2 header bar
\definecolor{clr3hdr}{RGB}{112,48,160}      % Cluster 3 header bar

\begin{landscape}
\setlength{\tabcolsep}{4pt}
\renewcommand{\arraystretch}{1.25}

\begin{longtable}{%
|>{\footnotesize\raggedright\arraybackslash}p{3.4cm}%
|>{\footnotesize\raggedright\arraybackslash}p{5.5cm}%
|>{\footnotesize\raggedright\arraybackslash}p{5.5cm}%
|>{\footnotesize\raggedright\arraybackslash}p{5.5cm}|}

\caption{Progression Matrix --- Priority Cluster~1: Sovereign Trusted Infrastructure (P1: Hardware RoT, P2: PQC Agility, P3: Lightweight Crypto \& PhySec, P8: Confidential Computing, P9: Island Mode). Each phase is assessed across four analytical dimensions. Critical Enabler rows state the single named gate condition controlling the transition into the following phase.}
\label{tab:cluster1_matrix}\\
\hline
\rowcolor{dimHdr}
\multicolumn{4}{|>{\footnotesize\bfseries\centering\arraybackslash}p{21.5cm}|}{%
  \textcolor{white}{\colorbox{clr1hdr}{\parbox{21.0cm}{\centering
  \textbf{Priority Cluster 1 --- Sovereign Trusted Infrastructure}
  \quad|\quad Sub-priorities: P1, P2, P3, P8, P9}}}} \\
\hline
\rowcolor{dimHdr}
\textbf{Analytical Dimension} &
\textbf{Phase 1: Foundation \& Compliance (2025--2030)} &
\textbf{Phase 2: Applied \& Industrial R\&D alongside Federation \& Standardisation (2030--2035)} &
\textbf{Phase 3: Autonomous Leadership \& Frontier Innovation (2035--2040)} \\
\hline
\endfirsthead

\hline
\rowcolor{dimHdr}
\textbf{Analytical Dimension} &
\textbf{Phase 1 (2025--2030)} &
\textbf{Phase 2 (2030--2035)} &
\textbf{Phase 3 (2035--2040)} \\
\hline
\endhead

\hline\multicolumn{4}{r}{\footnotesize\textit{Continued\ldots}}\\
\endfoot
\hline
\endlastfoot

% ---- R&I ROW ----
\rowcolor{dimRI}
\cellhead{(1) R\&I Maturity}
&
RISC-V secure element prototypes validated against EUCC evaluation criteria under EU Chips Act pilot funding; NIST FIPS~203/204/205 algorithms optimised for 8/16-bit microcontrollers within 32\,KB~RAM constraints; RF fingerprinting CSI/CFO device authentication validated in controlled mMTC environments; RISC-V Keystone TEE attestation protocols tested against heterogeneous edge fleet configurations; soft-state Island Mode credential caching validated in 5G disconnection simulation environments.
The cluster ODP is Applied R\&D --- the lowest common denominator driven by P2, P3, and P8; P1 and P9 have reached Fragmented Piloting and are not undeployed.
\kpistrip{Transition criterion: at least one functional prototype per sub-priority validated against the minimum EU compliance baseline (NIS2, CRA, or AI Act as applicable).}
&
EU-manufactured RISC-V secure elements complete EUCC~`High' evaluation and enter industrial production; hardware-optimised ML-KEM and ML-DSA validated on EU silicon for constrained-device targets; RISC-V Keystone TEEs deployed in multi-provider EU Data Spaces with continuous remote attestation; FHE performance benchmarks achieve operational viability for defined AI pipeline workloads; Island Mode swarm coordination protocols validated in multi-operator cross-border corridor deployments.
Advanced applied R\&D continues in parallel on zero-energy PhySec, neuromorphic secure compute, and FHE generalisation.
\kpistrip{Transition criterion: production-ready EUCC-evaluated reference implementation for P1; validated performance benchmarks on EU silicon for P2, P3, P8; 72-hour Island Mode survivability demonstrated operationally for P9.}
&
EU research groups set the international frontier agenda in lightweight PQC, confidential computing, and physical-layer security without adopting externally defined research directions. Second-generation PQC hardware co-designed with EU secure elements; neuromorphic secure edge compute integrated with zero-energy PhySec; quantum-native TEE architectures for post-CMOS environments; bio-inspired Island Mode swarm protocols achieving autonomous cyber-physical recovery at continental scale.
\kpistrip{Transition criterion: the capability is included as a mandatory baseline requirement in at least one EU-wide standard or procurement framework with no waiver provisions for new deployments.}
\\
\hline

% ---- POLICY ROW ----
\rowcolor{dimPol}
\cellhead{(2) Policy, Regulatory \& Standards}
&
CRA Articles~10 and~11 full enforcement from October~2027 establishes the mandatory market condition for hardware security certification investment. NIS2 Article~21(2)(e) cryptographic policy obligations active from October~2024. NIST FIPS~203/204/205 finalized August~2024 providing algorithm standard basis; ENISA PQC transition guidelines in publication; EUCC evaluation methodology available but no EU-manufactured secure element has completed High assurance evaluation; 3GPP~TS~33.501 Release~18 provides 5G security architecture for Island Mode integration. DORA Article~9 ICT resilience testing requirements apply to P9; AI Act Article~9 risk management obligations apply to TEE-attested autonomous workloads.
\kpistrip{Critical policy gap: no EU-specific ratified standard for RISC-V secure element security evaluation existed at the time of writing; this must be addressed in Phase~1 CEN/CENELEC JTC~13 work programme.}
&
EUCC~`High' certification of the first EU-manufactured RISC-V secure element issued --- the single most critical policy gate condition for this cluster. CEN/CENELEC hardware security standards for RISC-V secure elements ratified; ETSI~TS~103~744 hybrid key exchange profiles adopted as mandatory profiles in EU critical infrastructure procurement; EUCS~`High+' certification operational for TEE-backed cloud services; 3GPP Release~19/20 specifies PhySec and ISAC security requirements for 6G non-terrestrial network integration.
\kpistrip{Critical policy gate: EUCC `High' certification of the first EU-manufactured RISC-V secure element is the binding gate condition controlling the Phase~2 SSL transition for P1, P2, P3, and P8 simultaneously.}
&
EUCC~`High' mandatory baseline for all new EU ICT products in regulated sectors; EU PQC profiles referenced in ISO/IEC international standards, establishing EU as a standards-setting jurisdiction rather than an adopter; continuous automated EUCS/EUCC certification operational at continental scale without manual audit cycles; AI Act Article~9 conformity assessment for hardware-backed autonomous security AI fully operational and referenced internationally; DORA Article~11 recovery time objectives achieved autonomously by Island Mode architecture.
\kpistrip{End-state governance criterion: at least one EU standard in this cluster's capability domain adopted as or formally referenced in an ISO/IEC or ITU-T international standard.}
\\
\hline

% ---- OPERATIONAL ROW ----
\rowcolor{dimOps}
\cellhead{(3) Operational Deployment \& Adoption}
&
Intel TDX and AMD SEV-SNP dominate all production TEE deployments in EU cloud infrastructure; RISC-V secure elements at prototype stage in 3--5 EU-funded research consortia; edge nodes universally operate with hard cloud dependencies and no validated autonomous fallback in production; hybrid PQC pilots operate in cloud research environments but not in IoT field deployments; Island Mode soft-state protocols validated in 5G research corridors.
Deployment actors are limited to large institutional early adopters (national cybersecurity agencies, telecom operators in EU corridor pilots, EU-funded research consortia); SME adoption is absent due to tooling immaturity.
\kpistrip{Transition criterion: documented operational deployment across at least two EU critical infrastructure sectors or Member State jurisdictions with published findings.}
&
EU RISC-V secure elements deployed in energy and telecommunications critical infrastructure; hybrid ML-KEM plus ECDH-P384 key encapsulation operational in Industrial IoT corridor deployments across 3+ Member States; multi-party confidential computing pipelines active in 3+ EU Data Spaces with validation compliance verification; 72-hour Island Mode survivability validated operationally in 5G cross-border corridor deployments without WAN connectivity; SME-accessible crypto-agility OTA tooling available via Digital Europe Programme.
\kpistrip{Transition criterion: deployment across five or more Member States or three or more EU critical infrastructure sectors, with documented SME participation in at least one deployment context.}
&
At least 60\% of EU critical infrastructure uses EU-designed EUCC~`High' certified secure elements as default procurement; all new EU IoT deployments ship with FIPS~203/204/205-compliant PQC without bespoke integration overhead; FHE operationalised end-to-end for privacy-preserving AI pipelines from IoT to HPC; bio-inspired Island Mode swarms achieve DORA-mandated recovery objectives autonomously; Security-by-Design is structurally indistinguishable from product functionality across the EU ICT market.
\kpistrip{End-state deployment criterion: at least 80\% of new ICT deployments include this capability as a standard component without dedicated integration effort; the compliance distinction between ``secure'' and ``standard'' products ceases to be meaningful.}
\\
\hline

% ---- MARKET/SSL ROW ----
\rowcolor{dimMkt}
\cellhead{(4) Market Uptake \& Strategic Sovereignty}
\sslbadge{sslredbg}{sslredtx}{SSL: High Foreign Dependency}
&
Non-EU TPMs hold $>$90\% of EU critical infrastructure secure element market; EU Chips Act fabrication pilot lines initiated but first certified secure elements are 3--5 years from production; no EU-manufactured EUCC~`High' certified component available in any product category; PQC implementation ecosystem dominated by US/UK cryptography research outputs; Island Mode tooling absent from EU commercial product portfolios.
\kpistrip{SSL baseline: High Foreign Dependency across all five sub-priorities. Market transition criterion: at least one EU-designed solution achieves a verifiable market entry event (product launch, procurement award, or EUCC/EUCS certification).}
&
\sslbadge{sslamberbg}{sslambtx}{SSL: Emerging EU Ecosystem}
EU alternatives hold 15--30\% market share in new critical infrastructure secure element procurement; EU Chips Act fabrication ramping toward 50\% self-sufficiency target; first commercial EU TEE products competing at product level with Intel TDX/AMD SEV-SNP; foreign dependency reduced from structural to situational in regulated-sector procurement; EU industrial actors present in sovereign hardware procurement across 5+ Member States.
\kpistrip{Market transition criterion: EU-designed solutions hold at least 30\% of new critical infrastructure procurement market share, or foreign dependency formally reduced below the threshold identified at baseline.}
&
\sslbadge{ssltealbg}{ssltealtx}{SSL: Full Sovereign Autonomy}
EU Chips Act sovereign fabrication meets $\geq$60\% demand for security-critical silicon across all product categories; EU TEE, secure element, and PQC solutions commercially exported to partner jurisdictions; zero structural foreign dependency in new critical infrastructure deployments; EU sets technical terms of competition in secure silicon globally through EUCC certification requirements and procurement weight; foreign legacy deployments under active replacement programmes.
\kpistrip{End-state SSL criterion: EU-designed solutions hold at least 60\% of new critical infrastructure procurement; foreign legacy under binding replacement; EU sovereignty is self-sustaining, not programme-managed.}
\\
\hline

% ---- GATE ROW ----
\rowcolor{dimGate}
\textit{\textbf{Critical Enabler --- Phase~1 $\rightarrow$ Phase~2}}
&
\multicolumn{1}{>{\footnotesize\raggedright\arraybackslash\columncolor{dimGate}}p{5.5cm}|}{%
\textbf{CRA Articles~10 and~11 enforcement (October~2027)} creating mandatory market demand for hardware security certification and establishing the commercial conditions that justify EU Chips Act fabrication investment and EUCC~`High' evaluation resource allocation.}
&
\multicolumn{1}{>{\footnotesize\raggedright\arraybackslash\columncolor{dimGate}}p{5.5cm}|}{%
\textbf{EUCC~`High' certification of the first EU-manufactured RISC-V secure element} constituting the gate condition that transitions the cluster from research-grade to market-grade sovereign hardware and unlocks EU Chips Act fabrication scale-up investment justification.}
&
\multicolumn{1}{>{\footnotesize\raggedright\arraybackslash\columncolor{dimGate}}p{5.5cm}|}{%
\textbf{EU Chips Act sovereign EUV-capable fabrication achieving $\geq$60\% market self-sufficiency in security-critical silicon} removing the last structural foreign dependency in the infrastructure trust chain and enabling Full Sovereign Autonomy as a self-sustaining market condition.}
\\
\hline

% ---- SUB-PRIORITY DECOMPOSITION ----
\rowcolor{dimHdr}
\multicolumn{4}{|>{\footnotesize\itshape\raggedright\arraybackslash}p{21.5cm}|}{%
\textbf{Sub-Priority Decomposition} --- individual maturity tracking per constituent priority across phases.}
\\
\hline

\rowcolor{dimSub}
\textit{P1: Hardware Roots of Trust \& EUCC `High'}
& ODP: Fragmented Piloting \newline SSL: High Foreign Dependency
& ODP: Standardised Integration \newline SSL: Emerging EU Ecosystem
& ODP: Ubiquitous Operation \newline SSL: Full Sovereign Autonomy \\
\hline

\rowcolor{dimSub}
\textit{P2: PQC Agility for Edge/IoT}
& ODP: Applied R\&D \newline SSL: High Foreign Dependency
& ODP: Fragmented Piloting \newline SSL: Emerging EU Ecosystem
& ODP: Ubiquitous Operation \newline SSL: Full Sovereign Autonomy \\
\hline

\rowcolor{dimSub}
\textit{P3: Lightweight Crypto \& Physical-Layer Security}
& ODP: Applied R\&D \newline SSL: High Foreign Dependency
& ODP: Standardised Integration \newline SSL: Emerging EU Ecosystem
& ODP: Ubiquitous Operation \newline SSL: Full Sovereign Autonomy \\
\hline

\rowcolor{dimSub}
\textit{P8: Confidential Computing \& Privacy-Preserving AI}
& ODP: Applied R\&D \newline SSL: High Foreign Dependency
& ODP: Fragmented Piloting \newline SSL: Emerging EU Ecosystem
& ODP: Ubiquitous Operation \newline SSL: Full Sovereign Autonomy \\
\hline

\rowcolor{dimSub}
\textit{P9: Resilient-by-Design \& Autonomous Island Mode}
& ODP: Fragmented Piloting \newline SSL: High Foreign Dependency
& ODP: Standardised Integration \newline SSL: Contained Local Survivability
& ODP: Ubiquitous Operation \newline SSL: Systemic EU Resilience \\
\hline

\end{longtable}

\end{landscape}

\newpage

\section*{Appendix: Cluster Matrix Notes}

The priority matrices presented above organize strategic priorities across three phases (2025--2030, 2030--2035, 2035--2040) and four analytical dimensions: R\&D Maturity, Policy/Regulatory/Standards, Operational Deployment \& Adoption, and Market Uptake \& Strategic Sovereignty.

Critical Enabler rows identify the gate conditions controlling phase transitions. Each condition must be met to justify investment and resource allocation in the following phase.

\end{document}
